\section{Reviewer A comments}
{\it 
``The manuscript presents a new method to improve angular uncertainty in antineutrino direction reconstruction for segmented detectors. The authors apply a previously published algorithm to simulated data, optimize it for segmented detector geometries, and ultimately compare the results with observed data. Compared to conventional methods, the proposed approach aims to optimize for low-count limitations and provides a more robust way to handle error estimation.

The technical details are sound, and the results are clearly presented across various conditions, such as different segment sizes. The authors also provide a helpful elaboration of the assumptions underlying this work at the end of page 5. However, the manuscript would benefit from further clarification regarding the potential impact of these factors on the angular uncertainty estimation and the overall reconstruction. Given that these assumptions may have non-trivial effects, providing additional detail would significantly help to strengthen the robustness of this work.
''
}

\section{Response}

Dear Reviewer A:

Thank you for your comments and helpful feedback. We appreciate the positive assessment and your suggestions for improving the work. To that end, our response to your comments addresses the assumptions on page 5 listed below:

\begin{itemize}
    \item impact of 1/\(r^2\) assumption
    \item impact of angular spread of source
    \item impact of escaping neutrons
\end{itemize}

The following statement has been added to the conclusion highlighting a planned use case study in which the impacts of these assumptions are relevant and will be studied:

``Applying this method to different deployment scenarios, such as small modular reactors, will investigate some of the underlying assumptions, such as neglecting \(\frac{1}{r^2}\) dependence and non-zero angular size of nearby reactors are all being considered for future work.''

Sincerely,

Brian Crow
Department of Physics and Astronomy, University of Hawai'i at M\={a}noa

on behalf of the authors